
\documentclass{article}
\usepackage{colortbl}
\usepackage{makecell}
\usepackage{multirow}
\usepackage{supertabular}

\begin{document}

\newcounter{utterance}

\centering \large Interaction Transcript for game `function\_detective', experiment `list\_passive\_examples', episode 115 with gpt{-}5{-}nano.
\vspace{24pt}

{ \footnotesize  \setcounter{utterance}{1}
\setlength{\tabcolsep}{0pt}
\begin{supertabular}{c@{$\;$}|p{.15\linewidth}@{}p{.15\linewidth}p{.15\linewidth}p{.15\linewidth}p{.15\linewidth}p{.15\linewidth}}
    \# & \multicolumn{2}{c}{Player} && \multicolumn{2}{c}{Game Master} \\
    \hline

    \theutterance \stepcounter{utterance}  
    & & & \multicolumn{4}{p{0.6\linewidth}}{
        \cellcolor[rgb]{0.9,0.9,0.9}{
            \makecell[{{p{\linewidth}}}]{
                \texttt{\tiny{[P1$\langle$GM]}}
                \texttt{I have a hidden Python function for you to figure out.} \\
\\ 
\texttt{Your job is to infer what it does and then submit a correct Python implementation.} \\
\\ 
\texttt{Function signature:} \\
\texttt{(items: List{[}int{]}) {-}> List{[}int{]}} \\
\\ 
\texttt{You have 15 turns total.} \\
\texttt{Setup: Passive Examples} \\
\\ 
\texttt{I control which valid input{-}output examples you see. You can request them one at a time.} \\
\\ 
\texttt{You must make your best SOLVE attempt no later than the second{-}to{-}last turn.} \\
\texttt{Do not wait until the very end.} \\
\\ 
\texttt{You can use these actions:} \\
\texttt{{-} NEXT:} \\
\texttt{{-} SOLVE: ```python ... ```} \\
\\ 
\texttt{Final answer format:} \\
\texttt{SOLVE: ```python} \\
\texttt{def solution(<value1>):} \\
\texttt{    ...} \\
\texttt{```} \\
\\ 
\texttt{Do not add import statements in your guess.} \\
\texttt{Anything you might need is already available, so just write the function body and return the answer.} \\
\\ 
\texttt{To receive the next valid example, write exactly `NEXT:`.} \\
\texttt{Each reply is a valid pair (x, f(x)).} \\
\\ 
\texttt{I will reply with OUTPUT: lines only.} \\
\texttt{You must never write OUTPUT: yourself.} \\
\\ 
\texttt{Preloaded examples:} \\
\texttt{1. ({[}5, {-}5, 5{]}, {[}5, {-}5{]})} \\
\\ 
\texttt{Scoring:} \\
\texttt{{-} correctness is required for any non{-}zero score} \\
\texttt{{-} in interactive modes, solving with fewer information requests is better} \\
\\ 
\texttt{Input formatting reminders:} \\
\texttt{{-} strings must be quoted} \\
\texttt{{-} booleans must be True/False} \\
\texttt{{-} use commas between multiple values} \\
\texttt{{-} do not write any extra action tags in the same message} \\
\\ 
\texttt{Format example:} \\
\texttt{You write:} \\
\texttt{NEXT:} \\
\texttt{I answer:} \\
\texttt{OUTPUT: ({[}5, {-}5, 5{]}, {[}5, {-}5{]})} \\
\\ 
\texttt{Output only one valid action per turn.} \\
            }
        }
    }
     \\ \\

    \theutterance \stepcounter{utterance}  
    & \multicolumn{4}{p{0.6\linewidth}}{
        \cellcolor[rgb]{0.561,0.737,0.561}{
            \makecell[{{p{\linewidth}}}]{
                \texttt{\tiny{[P1$\rangle$GM]}}
                \texttt{SOLVE: ```python} \\
\texttt{def solution(items):} \\
\texttt{    \# Infer from example: input {[}5, {-}5, 5{]} {-}> output {[}5, {-}5{]}} \\
\texttt{    \# It seems to drop the last occurrence of the maximum absolute value? } \\
\texttt{    \# Here absolute values: {[}5,5,5{]}, all equal. It kept first two elements and dropped the last.} \\
\texttt{    \# Another plausible rule: remove trailing element with same sign as first? first is 5 positive; last is 5 positive so drop.} \\
\texttt{    \# Generalize: drop the last element if it has the same sign as the first element.} \\
\texttt{    if not items:} \\
\texttt{        return {[}{]}} \\
\texttt{    first\_sign = 1 if items{[}0{]} > 0 else ({-}1 if items{[}0{]} < 0 else 0)} \\
\texttt{    if first\_sign == 0:} \\
\texttt{        return items{[}:{-}1{]}} \\
\texttt{    if items{[}{-}1{]} * first\_sign > 0:} \\
\texttt{        return items{[}:{-}1{]}} \\
\texttt{    return items} \\
\texttt{```} \\
            }
        }
    }
    & & \\ \\

    \theutterance \stepcounter{utterance}  
    & & & \multicolumn{2}{p{0.3\linewidth}}{
        \cellcolor[rgb]{0.9,0.9,0.9}{
            \makecell[{{p{\linewidth}}}]{
                \texttt{\tiny{[GM$|$GM]}}
                \texttt{Game Verdict: LOSS} \\
            }
        }
    }
    & & \\ \\

    \theutterance \stepcounter{utterance}  
    & & & \multicolumn{2}{p{0.3\linewidth}}{
        \cellcolor[rgb]{0.9,0.9,0.9}{
            \makecell[{{p{\linewidth}}}]{
                \texttt{\tiny{[GM$|$GM]}}
                \texttt{Incorrect.} \\
\texttt{Failed Cases:} \\
\texttt{Input: {[}{[}0{]}{]} | Expected: {[}0{]} | Got: {[}{]}} \\
\texttt{Input: {[}{[}3, 2, 1{]}{]} | Expected: {[}3, 2, 1{]} | Got: {[}3, 2{]}} \\
\texttt{Input: {[}{[}1, 2, 3{]}{]} | Expected: {[}1, 2, 3{]} | Got: {[}1, 2{]}} \\
            }
        }
    }
    & & \\ \\

    \theutterance \stepcounter{utterance}  
    & & & \multicolumn{2}{p{0.3\linewidth}}{
        \cellcolor[rgb]{0.9,0.9,0.9}{
            \makecell[{{p{\linewidth}}}]{
                \texttt{\tiny{[GM$|$GM]}}
                \texttt{def rmv\_duplicates(items: List{[}int{]}) {-}> List{[}int{]}:} \\
\texttt{    seen = set()} \\
\texttt{    out: List{[}int{]} = {[}{]}} \\
\texttt{    for item in items:} \\
\texttt{        if item not in seen:} \\
\texttt{            seen.add(item)} \\
\texttt{            out.append(item)} \\
\texttt{    return out} \\
            }
        }
    }
    & & \\ \\

    \theutterance \stepcounter{utterance}  
    & & & \multicolumn{2}{p{0.3\linewidth}}{
        \cellcolor[rgb]{0.9,0.9,0.9}{
            \makecell[{{p{\linewidth}}}]{
                \texttt{\tiny{[GM$|$GM]}}
                \texttt{def solution(items):} \\
\texttt{    \# Infer from example: input {[}5, {-}5, 5{]} {-}> output {[}5, {-}5{]}} \\
\texttt{    \# It seems to drop the last occurrence of the maximum absolute value? } \\
\texttt{    \# Here absolute values: {[}5,5,5{]}, all equal. It kept first two elements and dropped the last.} \\
\texttt{    \# Another plausible rule: remove trailing element with same sign as first? first is 5 positive; last is 5 positive so drop.} \\
\texttt{    \# Generalize: drop the last element if it has the same sign as the first element.} \\
\texttt{    if not items:} \\
\texttt{        return {[}{]}} \\
\texttt{    first\_sign = 1 if items{[}0{]} > 0 else ({-}1 if items{[}0{]} < 0 else 0)} \\
\texttt{    if first\_sign == 0:} \\
\texttt{        return items{[}:{-}1{]}} \\
\texttt{    if items{[}{-}1{]} * first\_sign > 0:} \\
\texttt{        return items{[}:{-}1{]}} \\
\texttt{    return items} \\
            }
        }
    }
    & & \\ \\

    \theutterance \stepcounter{utterance}  
    & & & \multicolumn{2}{p{0.3\linewidth}}{
        \cellcolor[rgb]{0.9,0.9,0.9}{
            \makecell[{{p{\linewidth}}}]{
                \texttt{\tiny{[GM$|$GM]}}
                \texttt{Mode: passive\_examples} \\
\texttt{Accuracy: 0.380} \\
\texttt{Information requests: 0/15} \\
\texttt{Revealed passive examples: 1} \\
\texttt{Efficiency: 1.000} \\
\texttt{Slack: 15} \\
\texttt{Observed{-}pair consistency: 1.000} \\
\texttt{Parse errors: 0} \\
\texttt{Runtime errors: 0} \\
            }
        }
    }
    & & \\ \\

\end{supertabular}
}

\end{document}
